% i18n_qa_he_paper.tex — Hebrew computer science paper on algorithm complexity
% Intentional errors: missing geresh for abbreviations, RTL/LTR boundary
%   issues, missing bidi package
\documentclass[12pt,a4paper]{article}

\usepackage[english,hebrew]{babel}
\usepackage[utf8]{inputenc}
\usepackage[T1]{fontenc}
\usepackage{amsmath}
\usepackage{amssymb}
\usepackage{amsthm}
\usepackage{graphicx}
\usepackage{hyperref}
\usepackage[margin=2.5cm]{geometry}
\usepackage{booktabs}
\usepackage{natbib}
\usepackage{enumitem}
\usepackage{algorithm}
\usepackage{algpseudocode}
\usepackage{microtype}
\usepackage{xcolor}
\usepackage{float}
\usepackage{mathrsfs}
% INTENTIONAL ERROR: missing \usepackage{bidi} — needed for proper RTL

\theoremstyle{plain}
\newtheorem{theorem}{משפט}[section]
\newtheorem{lemma}[theorem]{למה}
\newtheorem{proposition}[theorem]{טענה}
\newtheorem{corollary}[theorem]{מסקנה}

\theoremstyle{definition}
\newtheorem{definition}[theorem]{הגדרה}
\newtheorem{example}[theorem]{דוגמה}

\theoremstyle{remark}
\newtheorem{remark}[theorem]{הערה}

\title{חסמים תחתונים לסיבוכיות אלגוריתמים\\
למיון וחיפוש במבני נתונים דינמיים}
\author{%
  % INTENTIONAL ERROR: missing geresh for abbreviation פרופ (should be פרופ׳)
  פרופ דוד כהן\thanks{המחלקה למדעי המחשב, האוניברסיטה העברית בירושלים.}
  \and
  % INTENTIONAL ERROR: missing geresh for דר (should be ד״ר)
  דר רחל לוי\thanks{מכון ויצמן למדע, רחובות.}
}
\date{6 באפריל 2026}

\begin{document}

\maketitle

\begin{abstract}
במאמר זה אנו מוכיחים חסמים תחתונים חדשים לסיבוכיות הזמן של אלגוריתמים
למיון והכנסה במבני נתונים דינמיים. אנו מראים כי כל מבנה נתונים דינמי
התומך בפעולות הכנסה, מחיקה ושאילתת קודם במודל השוואות חייב לבצע
לפחות $\Omega(\log n / \log \log n)$ פעולות לכל עדכון. תוצאה זו משפרת
את החסם התחתון הקלאסי של $\Omega(\log n)$ של פרדמן ושהם. בנוסף, אנו
מציגים מבנה נתונים חדש המשיג חסם זה עד גורם קבוע, מה שמוכיח כי
החסם התחתון שלנו הדוק. השיטה מבוססת על תורת המידע ועל טכניקת כרונוגרמות.
\end{abstract}

\tableofcontents

%% ---------------------------------------------------------------
\section{מבוא}
\label{sec:intro}

סיבוכיות חישובית היא ענף מרכזי במדעי המחשב התיאורטיים העוסק
בחקר המשאבים הנדרשים לפתרון בעיות חישוביות. במאמר זה אנו
מתמקדים בסיבוכיות מבני נתונים דינמיים --- מבנים התומכים
בעדכונים ובשאילתות באופן מקוון.

% INTENTIONAL ERROR: RTL/LTR boundary issue — English in Hebrew context without proper isolation
בעיית המיון הדינמי, הידועה גם כ-dynamic sorting problem, היא
אחת הבעיות היסודיות ביותר. במודל ההשוואות, כל אלגוריתם מיון
סטטי דורש $\Omega(n \log n)$ השוואות כדי למיין $n$ איברים,
כפי שהוכח על ידי עץ ההחלטות\footnote{%
ראו \citet{knuth1998}, כרך~3, פרק~5.3, למבוא מקיף.}.

\subsection{הגדרות ותוצאות קודמות}

\begin{definition}
\label{def:dynamic_ds}
מבנה נתונים דינמי למיון הוא מבנה $D$ התומך בפעולות הבאות:
\begin{enumerate}[label=(\alph*)]
  \item $\textsc{Insert}(x)$: הכנסת איבר $x$ למבנה;
  \item $\textsc{Delete}(x)$: מחיקת איבר $x$ מהמבנה;
  \item $\textsc{Predecessor}(x)$: מציאת האיבר הגדול ביותר $y \in D$ כך ש-$y \leq x$.
\end{enumerate}
\end{definition}

% INTENTIONAL ERROR: missing geresh for abbreviation עמ (should be עמ׳)
התוצאה הקלאסית של פרדמן ושהם (1989, עמ 42--58) קובעת כי
במודל עץ החלטות, כל מבנה נתונים דינמי התומך בהכנסה ובשאילתת
קודם דורש זמן $\Omega(\log n)$ לכל פעולה, כאשר $n$ הוא מספר
האיברים במבנה~\citep{fredman1989}.

\subsection{תרומות המאמר}

תרומותינו העיקריות הן:

\begin{theorem}[חסם תחתון עיקרי]
\label{thm:main}
כל מבנה נתונים דינמי במודל ההשוואות התומך בפעולות
$\textsc{Insert}$, $\textsc{Delete}$ ו-$\textsc{Predecessor}$
חייב לבצע לפחות
\begin{equation}
\label{eq:lower_bound}
T(n) = \Omega\!\left(\frac{\log n}{\log \log n}\right)
\end{equation}
פעולות ממוצעות (amortized) לכל עדכון.
\end{theorem}

\begin{theorem}[חסם עליון תואם]
\label{thm:upper}
קיים מבנה נתונים דינמי במודל ההשוואות המשיג
\begin{equation}
\label{eq:upper_bound}
T(n) = O\!\left(\frac{\log n}{\log \log n}\right)
\end{equation}
פעולות לכל עדכון ושאילתה.
\end{theorem}

מבנה הנייר הוא כדלקמן: סעיף~\ref{sec:model} מגדיר את המודל
הפורמלי, סעיף~\ref{sec:lower} מוכיח את החסם התחתון,
סעיף~\ref{sec:upper} מתאר את מבנה הנתונים, וסעיף~\ref{sec:apps}
דן ביישומים.

%% ---------------------------------------------------------------
\section{המודל הפורמלי}
\label{sec:model}

\subsection{מודל ההשוואות}

במודל ההשוואות, האלגוריתם יכול לגשת לנתונים רק באמצעות
השוואות בין זוגות איברים. באופן פורמלי:

\begin{definition}
\label{def:comparison_model}
אלגוריתם במודל ההשוואות מוגדר כעץ החלטות בינארי $\mathcal{T}$
בו כל צומת פנימי מכיל השוואה $x_i \stackrel{?}{\leq} x_j$
וכל עלה מציין פלט. הסיבוכיות היא עומק העץ.
\end{definition}

\subsection{מודל הכרונוגרמה}

% INTENTIONAL ERROR: RTL/LTR boundary — English terms without isolation
הכרונוגרמה (chronogram technique) היא שיטה להוכחת חסמים
תחתונים שפותחה על ידי פרדמן ושהם. הרעיון המרכזי הוא
לשקול את ה-information flow בין פעולות עדכון לשאילתות.

\begin{definition}
\label{def:chronogram}
כרונוגרמה של סדרת פעולות $\sigma = (o_1, o_2, \ldots, o_m)$
על מבנה נתונים $D$ היא חלוקה של צעדי החישוב לשכבות
$L_0, L_1, \ldots, L_t$ כך ש-$L_i$ מכילה את הצעדים
שבוצעו בזמנים $[2^i, 2^{i+1})$ לאחר העדכון האחרון.
\end{definition}

\subsection{אנטרופיה ומידע}

הוכחת החסם התחתון משתמשת בתורת המידע. נזכיר הגדרות בסיסיות:

\begin{equation}
\label{eq:entropy}
H(X) = -\sum_{x} p(x) \log_2 p(x),
\end{equation}
ו-המידע ההדדי:
\begin{equation}
\label{eq:mutual_info}
I(X; Y) = H(X) - H(X \mid Y) = H(Y) - H(Y \mid X).
\end{equation}

לפי אי-שוויון פאנו\footnote{%
אי-שוויון פאנו קובע כי $H(X \mid Y) \leq h(P_e) + P_e \log(|\mathcal{X}| - 1)$,
ראו \citet{cover2006}.},
\begin{equation}
\label{eq:fano}
P_e \geq \frac{H(X \mid Y) - 1}{\log |\mathcal{X}|}.
\end{equation}

%% ---------------------------------------------------------------
\section{הוכחת החסם התחתון}
\label{sec:lower}

\subsection{בניית הקלט הקשה}

נבנה סדרת פעולות שתכריח כל אלגוריתם לבצע עבודה רבה.
נבחר $n$ איברים $x_1 < x_2 < \cdots < x_n$ ופרמוטציה
אקראית $\pi \in S_n$.

\begin{lemma}
\label{lem:info_content}
סדרת ההכנסות $\textsc{Insert}(x_{\pi(1)}), \textsc{Insert}(x_{\pi(2)}),
\ldots, \textsc{Insert}(x_{\pi(n)})$ מכילה $\log_2(n!) = \Theta(n \log n)$
סיביות מידע.
\end{lemma}

\begin{proof}
מספר הפרמוטציות הוא $n!$, לכן האנטרופיה היא $\log_2(n!)$.
לפי קירוב סטירלינג, $\log_2(n!) = n \log_2 n - n \log_2 e + O(\log n) = \Theta(n \log n)$.
\end{proof}

\subsection{ניתוח הכרונוגרמה}

\begin{lemma}
\label{lem:chronogram_layers}
בכרונוגרמה עם $t = \log n / \log \log n$ שכבות, כל שכבה $L_i$
חייבת להעביר לפחות
\[
\frac{n \log n}{t} = \Omega(n \log \log n)
\]
סיביות מידע.
\end{lemma}

\begin{proof}
סך המידע שיש להעביר הוא $\Theta(n \log n)$ סיביות
(למה~\ref{lem:info_content}). חלוקה שווה ל-$t$ שכבות נותנת
$n \log n / t$ סיביות לשכבה. עבור $t = \log n / \log \log n$:
\[
\frac{n \log n}{t} = \frac{n \log n}{\log n / \log \log n} = n \log \log n.
\]
\end{proof}

\begin{proof}[הוכחת משפט~\ref{thm:main}]
נשתמש בלמה~\ref{lem:chronogram_layers}. כל פעולת עדכון מבצעת
לכל היותר $T(n)$ השוואות, כאשר כל השוואה מעבירה סיבית אחת
של מידע. לכן, $n$ פעולות עדכון מעבירות לכל היותר $n \cdot T(n)$
סיביות. מכאן:
\begin{align}
\label{eq:main_proof}
n \cdot T(n) &\geq \frac{n \log n}{t} = n \log \log n, \notag \\
T(n) &\geq \log \log n = \frac{\log n}{\log n / \log \log n} = \frac{\log n}{\log n / \log \log n}.
\end{align}
עידון הניתוח באמצעות אי-שוויון פאנו~\eqref{eq:fano} מראה
שהחסם ההדוק הוא $\Omega(\log n / \log \log n)$.
\end{proof}

%% ---------------------------------------------------------------
\section{מבנה הנתונים העליון}
\label{sec:upper}

\subsection{תיאור המבנה}

מבנה הנתונים שלנו מבוסס על עצי חיפוש מאוזנים עם גורם
הסתעפות $B = \log^\epsilon n$ לערך קבוע $\epsilon > 0$.

\begin{definition}
\label{def:b_tree}
עץ-$B$ דינמי הוא עץ חיפוש מאוזן בו כל צומת פנימי מכיל
בין $B/2$ ל-$B$ מפתחות, וכל העלים נמצאים באותו עומק.
\end{definition}

העומק של עץ-$B$ עם $n$ מפתחות הוא:
\begin{equation}
\label{eq:depth}
h = O\!\left(\frac{\log n}{\log B}\right) = O\!\left(\frac{\log n}{\epsilon \log \log n}\right) = O\!\left(\frac{\log n}{\log \log n}\right).
\end{equation}

\subsection{אלגוריתם ההכנסה}

\begin{algorithm}[ht]
\caption{הכנסה לעץ-$B$ דינמי}
\label{alg:insert}
\begin{algorithmic}[1]
\Require מפתח $x$, עץ $T$
\Ensure עץ מעודכן $T'$
\State מצא את העלה $v$ המתאים ל-$x$ בעץ $T$
\State הכנס את $x$ לצומת $v$
\While{$v$ מכיל יותר מ-$B$ מפתחות}
  \State פצל את $v$ לשני צמתים
  \State העלה את המפתח האמצעי לצומת האב
  \State $v \gets \text{אב}(v)$
\EndWhile
\end{algorithmic}
\end{algorithm}

\begin{proposition}
\label{prop:insert_time}
אלגוריתם~\ref{alg:insert} מבצע $O(\log n / \log \log n)$
השוואות לכל הכנסה.
\end{proposition}

\begin{proof}
עומק העץ הוא $O(\log n / \log \log n)$ לפי~\eqref{eq:depth}.
בכל צומת, מציאת המיקום הנכון דורשת $O(\log B) = O(\log \log n)$
השוואות באמצעות חיפוש בינארי. סך ההשוואות:
\[
O\!\left(\frac{\log n}{\log \log n}\right) \cdot O(\log \log n) = O(\log n).
\]
אולם, באמצעות טכניקת fusion trees~\citep{fredman1993} ניתן
לחפש בכל צומת ב-$O(1)$ השוואות, ואז הסיבוכיות הכוללת היא
$O(\log n / \log \log n)$.
\end{proof}

%% ---------------------------------------------------------------
\section{יישומים}
\label{sec:apps}

\subsection{מיון דינמי}

\begin{table}[ht]
\centering
\caption{השוואת סיבוכיות מבני נתונים דינמיים}
\label{tab:comparison}
\begin{tabular}{lccc}
\toprule
מבנה נתונים & הכנסה & מחיקה & שאילתה \\
\midrule
עץ AVL & $O(\log n)$ & $O(\log n)$ & $O(\log n)$ \\
עץ אדום-שחור & $O(\log n)$ & $O(\log n)$ & $O(\log n)$ \\
Skip list & $O(\log n)$ & $O(\log n)$ & $O(\log n)$ \\
% INTENTIONAL ERROR: RTL/LTR boundary issue — parenthetical English
Fusion tree (פרדמן-וילארד) & $O(\log_w n)$ & $O(\log_w n)$ & $O(\log_w n)$ \\
המבנה שלנו & $O\!\left(\frac{\log n}{\log \log n}\right)$ & $O\!\left(\frac{\log n}{\log \log n}\right)$ & $O\!\left(\frac{\log n}{\log \log n}\right)$ \\
\bottomrule
\end{tabular}
\end{table}

טבלה~\ref{tab:comparison} משווה את מבנה הנתונים שלנו למבנים
קיימים. במודל ההשוואות, המבנה שלנו משיג את הסיבוכיות
האופטימלית~\citep{patrascu2008}.

\subsection{גרפים דינמיים}

מבנה הנתונים שלנו יכול לשמש כמרכיב בסיסי באלגוריתמים
לגרפים דינמיים. למשל, בעיית הקישוריות הדינמית דורשת
תחזוקת יער פורש בעת הכנסה ומחיקה של קשתות.

\begin{figure}[ht]
\centering
\fbox{\parbox{0.7\textwidth}{%
\centering
\textbf{מבנה העץ-$B$ הדינמי}\\[8pt]
\begin{tabular}{ccccc}
& & [25, 50] & & \\
& / & & \textbackslash & \\
{[5, 12, 18]} & & {[30, 38, 42]} & & {[55, 63, 78, 91]} \\
\end{tabular}\\[6pt]
\emph{דוגמה לעץ-$B$ עם $B = 4$. כל צומת מכיל 2--4 מפתחות.}
}}
\caption{עץ-$B$ דינמי עם גורם הסתעפות $B = 4$ המכיל 12 מפתחות.
העומק הוא 2, מה שתואם את החסם $O(\log n / \log B)$.}
\label{fig:btree}
\end{figure}

\subsection{יישום לגיאומטריה חישובית}

% INTENTIONAL ERROR: missing geresh for abbreviation ד (should be ד׳)
בגיאומטריה חישובית ב-$d$ ממדים (עבור $d$ קבוע), מבנה
הנתונים שלנו מאפשר תחזוקת מעטפת קמורה דינמית
בסיבוכיות $O(\log^{d-1} n / \log \log n)$ לכל עדכון.

\begin{figure}[ht]
\centering
\fbox{\parbox{0.7\textwidth}{%
\centering
\textbf{סיבוכיות כפונקציה של $n$}\\[8pt]
\begin{tabular}{cc}
$n$ & $\log n / \log \log n$ \\
\hline
$10^3$ & 3.01 \\
$10^6$ & 4.31 \\
$10^9$ & 5.17 \\
$10^{12}$ & 5.85 \\
$10^{15}$ & 6.39 \\
\end{tabular}\\[6pt]
\emph{הצמיחה של $\log n / \log \log n$ לעומת $\log n$}
}}
\caption{השוואה מספרית בין $\log n / \log \log n$ ל-$\log n$
עבור ערכי $n$ טיפוסיים. היתרון גדל עם גודל הקלט.}
\label{fig:growth}
\end{figure}

%% ---------------------------------------------------------------
\section{סיכום}
\label{sec:conclusion}

הוכחנו חסם תחתון חדש של $\Omega(\log n / \log \log n)$ לסיבוכיות
מבני נתונים דינמיים במודל ההשוואות (משפט~\ref{thm:main}).
הראינו שחסם זה הדוק על ידי הצגת מבנה נתונים התואם
(משפט~\ref{thm:upper}). התוצאה משפרת את החסם הקלאסי
$\Omega(\log n)$ ומספקת תמונה שלמה יותר של סיבוכיות
מבני נתונים דינמיים.

כיוונים עתידיים כוללים הכללה למודלים חזקים יותר כמו
מודל RAM ומודל cell probe, וכן יישום לבעיות ספציפיות
בגיאומטריה חישובית ובתורת הגרפים\footnote{%
ראו \citet{demaine2007} לסקירה מקיפה של חסמים תחתונים
למבני נתונים.}.

%% ---------------------------------------------------------------
\begin{thebibliography}{99}

\bibitem[Cover and Thomas(2006)]{cover2006}
Cover, T.\,M. and Thomas, J.\,A. (2006).
\emph{Elements of Information Theory}, 2nd ed.
Wiley-Interscience.

\bibitem[Demaine(2007)]{demaine2007}
Demaine, E.\,D. (2007).
Cache-oblivious algorithms and data structures.
In \emph{Lecture Notes from the EEF Summer School on Massive Data Sets}.

\bibitem[Fredman and Saks(1989)]{fredman1989}
Fredman, M.\,L. and Saks, M.\,E. (1989).
The cell probe complexity of dynamic data structures.
\emph{Proceedings of the 21st ACM STOC}, 345--354.

\bibitem[Fredman and Willard(1993)]{fredman1993}
Fredman, M.\,L. and Willard, D.\,E. (1993).
Surpassing the information theoretic bound with fusion trees.
\emph{Journal of Computer and System Sciences}, 47(3), 424--436.

\bibitem[Knuth(1998)]{knuth1998}
Knuth, D.\,E. (1998).
\emph{The Art of Computer Programming, Vol.~3: Sorting and Searching},
2nd ed. Addison-Wesley.

\bibitem[P\u{a}tra\c{s}cu(2008)]{patrascu2008}
P\u{a}tra\c{s}cu, M. (2008).
Lower bounds for 2-dimensional range counting.
\emph{Proceedings of the 39th ACM STOC}, 40--46.

\end{thebibliography}

\end{document}
